\documentclass[11pt]{article}
\usepackage[margin=1in]{geometry}
\usepackage{graphicx}
\usepackage{booktabs}
\usepackage{xcolor}
\usepackage{hyperref}
\usepackage{amsmath}
\usepackage{listings}
\usepackage{enumitem}
\usepackage{caption}
\usepackage{float}
\usepackage{longtable}
\usepackage{xurl}
\setlength{\emergencystretch}{3em}

\hypersetup{colorlinks=true, linkcolor=blue, urlcolor=blue, citecolor=blue}

\definecolor{codebg}{rgb}{0.96,0.96,0.96}
\lstset{
  basicstyle=\ttfamily\small,
  backgroundcolor=\color{codebg},
  frame=single,
  breaklines=true,
  columns=fullflexible,
  tabsize=2,
}

% A value that was configured but is not visible in any screenshot.
\newcommand{\fillin}{\textbf{\color{red!60!black}[fill in]}}

\newcommand{\shot}[3][\textwidth]{%
  \begin{figure}[H]\centering\includegraphics[width=#1]{figures/#2}\caption{#3}\label{fig:#2}\end{figure}}

\title{\Large\textbf{Lab Assignment Report: Building and Testing\\ Networks in Cisco Packet Tracer}}
\author{Sajjad Ahmed Shaaz \\ 002410501093 \\ BCSE-III A3}
\date{November 17, 2026}

\begin{document}
\maketitle

\section*{Basic Information}
\begin{itemize}[leftmargin=1.6em]
  \item \textbf{Subject:} CSE/PC/B/S/314 Computer Networks Lab
  \item \textbf{Assignment Number:} 6
  \item \textbf{Problem Statement:} Use Cisco Packet Tracer to build progressively larger
  networks: two hosts back-to-back, a hub LAN, a switch LAN (with ARP and MAC tables), the
  two LANs joined, two routed LANs with static routes, DHCP servers in those LANs, and finally
  a departmental LAN with DHCP, DNS, web and FTP servers under the domain
  \texttt{cse.myuniv.edu}.
  \item \textbf{Submission due:} 2nd lab after the Puja vacation (tentatively 02--06
  November 2026)
\end{itemize}

\section{Design}

\subsection{Purpose}
Assignment 5 observed a real network from the outside with Wireshark. This assignment
builds networks from the inside. Every device, cable, address, table and service is placed
and configured by hand, then tested with \texttt{ping}, \texttt{arp}, \texttt{ipconfig},
\texttt{show mac address-table}, a browser and an FTP client. The seven questions climb
the stack: Layer 1 (cables and hubs), Layer 2 (switching and ARP), Layer 3 (routing), and
finally application services (DHCP, DNS, HTTP, FTP).

\subsection{Environment}
\begin{table}[H]
\centering
\begin{tabular}{ll}
\toprule
\textbf{Item} & \textbf{Value} \\
\midrule
Simulator & Cisco Packet Tracer (Linux), Realtime mode \\
Saved topology & \texttt{assignment6.pkt} (all seven networks in one file) \\
End hosts & \texttt{PC-PT} (Windows-like command prompt, TTL 128) \\
Servers & \texttt{Server-PT} (DHCP, DNS, HTTP, FTP services) \\
Hub / generic switch & \texttt{Hub-PT}, \texttt{Switch-PT} (Q2--Q4) \\
Managed switch & Cisco \texttt{2950-24} (Q5--Q7) \\
Router & Cisco \texttt{1841} (Q5--Q6); see note below \\
\bottomrule
\end{tabular}
\caption{Devices used.}
\end{table}

\paragraph{Note on the router model.} The assignment asks for a ``Cisco 1818''. No such
model exists in Packet Tracer's router list (the closest are the 1841 and the 1941). The
\textbf{1841} was used. It has two FastEthernet ports, Fa0/0 for the LAN and Fa0/1 for
the router-to-router link, which is all that Q5 needs.

\subsection{Addressing Plan}
Table~\ref{tab:plan} summarises every network built. Per-device addresses are given in each
question's \emph{Topology and Addressing} subsection.

\begin{table}[H]
\centering
\small
\begin{tabular}{llp{4.3cm}l}
\toprule
\textbf{Q} & \textbf{Network} & \textbf{Subnet} & \textbf{Address assignment} \\
\midrule
1 & PC0 -- PC1 (cross-over) & 192.168.1.0/24 & static, .1--.2 \\
2 & LAN-A (hub) & 192.168.1.0/24 & static, .3--.5 \\
3 & LAN-B (switch) & 192.168.1.0/24 & static, .6--.8 \\
4 & LAN-A + LAN-B & 192.168.1.0/24 & unchanged (one broadcast domain) \\
5 & JU-Main / JU-SL / WAN & 192.168.148.0/24, .149.0/24, .150.0/24 & static, gateway \texttt{.1} \\
6 & JU-Main / JU-SL & same as Q5 & DHCP from a server in each LAN \\
7 & CSE & 192.168.151.0/24 & DHCP for hosts, static for servers \\
\bottomrule
\end{tabular}
\caption{Subnets used in the assignment.}
\label{tab:plan}
\end{table}

LAN-A and LAN-B deliberately share the subnet 192.168.1.0/24 (with no overlapping host
addresses). Q4 joins them at Layer 2 only, with no router. That works only if every host is
on the same IP subnet.

\subsection{Structure of Each Question}
Each question is written up in three parts:
\begin{enumerate}[leftmargin=1.8em]
  \item \textbf{Topology and Addressing}: a diagram of the network (drawn from the Packet
  Tracer workspace, with every address on it) and a table of devices.
  \item \textbf{Procedure}: what was placed, cabled and configured, and how.
  \item \textbf{Results}: screenshots of the tests, followed by observations.
\end{enumerate}
In the diagrams, a \textbf{solid line} is a copper straight-through cable and a
\textbf{dashed line} is a copper cross-over cable, matching Packet Tracer's own notation.

%==========================================================================
\section{Questions}

%--------------------------------------------------------------------------
\subsection{Q1: Two Hosts Back-to-Back with a Cross-over Cable}

\subsubsection*{Topology and Addressing}
\shot[0.75\textwidth]{topo_q1.png}{Q1 topology: two PCs joined directly by a cross-over cable.}

\begin{table}[H]
\centering
\begin{tabular}{llll}
\toprule
\textbf{Device} & \textbf{Interface} & \textbf{IPv4 address} & \textbf{Subnet mask} \\
\midrule
PC0 & FastEthernet0 & 192.168.1.1 & 255.255.255.0 \\
PC1 & FastEthernet0 & 192.168.1.2 & 255.255.255.0 \\
\bottomrule
\end{tabular}
\end{table}

\subsubsection*{Procedure}
\begin{enumerate}[leftmargin=1.8em]
  \item Placed two \texttt{PC-PT} end devices, PC0 and PC1.
  \item Selected \emph{Connections $\to$ Copper Cross-Over} and joined FastEthernet0 of
  PC0 to FastEthernet0 of PC1. Both link LEDs turned green.
  \item On each PC, \emph{Desktop $\to$ IP Configuration $\to$ Static}: entered the
  address above. Packet Tracer fills in the /24 mask automatically. No gateway is needed,
  because both hosts are on the same subnet.
  \item \emph{Desktop $\to$ Command Prompt}: \texttt{ping 192.168.1.2} from PC0, then
  \texttt{ping 192.168.1.1} from PC1.
\end{enumerate}

\subsubsection*{Results}
\shot{q1_ping_pc0.png}{PC0 pings PC1 (192.168.1.2): 4 sent, 4 received, 0\% loss.}
\shot{q1_ping_pc1.png}{PC1 pings PC0 (192.168.1.1): 4 sent, 4 received, 0\% loss.}

\paragraph{Observations.} Both directions succeed with RTT $<$1\,ms and TTL\,=\,128 (the
Windows default; no router decremented it). A cross-over cable is required here. A NIC
transmits on pins 1--2 and receives on pins 3--6, so connecting two NICs directly needs the
pairs swapped in the cable. A hub or switch normally does this swap internally.

%--------------------------------------------------------------------------
\subsection{Q2: LAN-A, Three Hosts on a Hub}

\subsubsection*{Topology and Addressing}
\shot[0.8\textwidth]{topo_q2.png}{Q2 topology: LAN-A, three PCs on a \texttt{Hub-PT}.}

\begin{table}[H]
\centering
\begin{tabular}{llll}
\toprule
\textbf{Device} & \textbf{IPv4 address / mask} & \textbf{MAC address} & \textbf{Connected to} \\
\midrule
PC2 & 192.168.1.3/24 & 00D0.FFA5.37EE & Hub (straight-through) \\
PC3 & 192.168.1.4/24 & 0090.2BE8.B750 & Hub (straight-through) \\
PC4 & 192.168.1.5/24 & 0001.C70A.0697 & Hub (straight-through) \\
\bottomrule
\end{tabular}
\caption{LAN-A hosts. MAC addresses are from the ARP tables recorded in Q4.}\label{tab:lana}
\end{table}

\subsubsection*{Procedure}
\begin{enumerate}[leftmargin=1.8em]
  \item Placed a \texttt{Hub-PT} (later relabelled ``LAN A'') and three PCs.
  \item Connected each PC to the hub with a \textbf{copper straight-through} cable (host to
  hub is a ``different device'' connection).
  \item Assigned static addresses 192.168.1.3--.5/24.
  \item Pinged every pair, in both directions: PC2$\to$PC3, PC2$\to$PC4, PC3$\to$PC2,
  PC3$\to$PC4, PC4$\to$PC2, PC4$\to$PC3.
\end{enumerate}

\subsubsection*{Results}
\shot{q2_ping_pc2.png}{PC2 pings PC3 (192.168.1.4) and PC4 (192.168.1.5). The very first
attempt shows one lost request, after which all replies arrive.}
\shot{q2_ping_all.png}{All pairs: PC2 (left) pings .4 and .5; PC4 (top right) pings .3
and .4; PC3 (bottom right) pings .3 and .5. Every run ends with 0\% loss.}

\paragraph{Observations.}
\begin{itemize}[leftmargin=1.6em]
  \item All six ordered pairs reach each other.
  \item The first-ever ping to an address may lose its first request (PC2$\to$.4 above;
  PC2 interrupted that run with Ctrl-C after one lost packet). The sender must first
  broadcast an ARP request, and Packet Tracer drops the ICMP packet that triggered the
  ARP rather than queuing it. Once the ARP reply is cached, every request succeeds.
  \item A hub is a Layer-1 repeater. Every frame, unicast or broadcast, is repeated out of
  every other port, so all three PCs share one collision domain and see each other's traffic.
\end{itemize}

%--------------------------------------------------------------------------
\subsection{Q3: LAN-B, Three Hosts on a Switch (ARP and MAC Tables)}

\subsubsection*{Topology and Addressing}
\shot[0.8\textwidth]{topo_q3.png}{Q3 topology: LAN-B, three PCs on a \texttt{Switch-PT}, with
the switch port of each host.}

\begin{table}[H]
\centering
\begin{tabular}{llll}
\toprule
\textbf{Device} & \textbf{IPv4 address / mask} & \textbf{MAC address} & \textbf{Switch port} \\
\midrule
PC5 & 192.168.1.6/24 & 0030.F223.90E0 & FastEthernet0/1 \\
PC6 & 192.168.1.7/24 & 0005.5ED3.1D1E & FastEthernet1/1 \\
PC7 & 192.168.1.8/24 & 0002.16AA.B4DC & FastEthernet2/1 \\
\bottomrule
\end{tabular}
\caption{LAN-B hosts. Ports are read from the switch's MAC table.}\label{tab:lanb}
\end{table}

\subsubsection*{Procedure}
\begin{enumerate}[leftmargin=1.8em]
  \item Placed a \texttt{Switch-PT} (later relabelled ``LAN B'') and PCs PC5--PC7, connected
  with straight-through cables, and assigned 192.168.1.6--.8/24.
  \item \textbf{Before any traffic:} ran \texttt{arp -a} on every PC and, in the switch
  CLI, \texttt{enable} followed by \texttt{show mac address-table}.
  \item Pinged every pair of hosts.
  \item \textbf{After the pings:} repeated \texttt{arp -a} on each PC and
  \texttt{show mac address-table} on the switch.
\end{enumerate}

\subsubsection*{Results: before pinging}
\shot{q3_arp_before.png}{\texttt{arp -a} on PC5, PC6 and PC7 before any traffic: ``No ARP
Entries Found'' on all three. (PC5 also ran \texttt{arp -d} to be sure the cache was empty.)}
\shot{q3_mac_before.png}{Switch CLI before any traffic. The first attempt,
\texttt{show mac-address table}, is rejected (wrong hyphen position). The correct
\texttt{show mac address-table} returns an empty table. The amber link lights show the
ports still coming up.}

\subsubsection*{Results: pinging}
\shot{q3_ping_pc5.png}{PC5 pings PC6 (192.168.1.7) and PC7 (192.168.1.8): 0\% loss.}
\shot{q3_mac_after_pc5.png}{MAC table after PC5's pings (left): the switch has learned
\texttt{0030.f223.90e0} (PC5) on Fa0/1 and \texttt{0002.16aa.b4dc} (PC7) on Fa2/1.}
\shot{q3_ping_all.png}{The remaining pairs: PC6 (top left) pings .6 and .8, PC5 (top right)
pings .7 and .8, PC7 (bottom right, starting from an empty ARP cache) pings .6 and .7. The
last \texttt{show mac address-table} now holds all three hosts.}

\subsubsection*{Results: after pinging}
\shot{q3_arp_after.png}{\texttt{arp -a} after all pairs were pinged. Each PC has exactly two
dynamic entries, one for each of the other two hosts.}

\begin{table}[H]
\centering
\small
\begin{tabular}{lll}
\toprule
\textbf{Host} & \textbf{ARP table before} & \textbf{ARP table after} \\
\midrule
PC5 (.6) & empty & .7 $\to$ 0005.5ed3.1d1e,\; .8 $\to$ 0002.16aa.b4dc \\
PC6 (.7) & empty & .6 $\to$ 0030.f223.90e0,\; .8 $\to$ 0002.16aa.b4dc \\
PC7 (.8) & empty & .6 $\to$ 0030.f223.90e0,\; .7 $\to$ 0005.5ed3.1d1e \\
\bottomrule
\end{tabular}
\vspace{0.6em}

\begin{tabular}{llll}
\toprule
\textbf{VLAN} & \textbf{MAC address} & \textbf{Type} & \textbf{Port} \\
\midrule
1 & 0002.16aa.b4dc (PC7) & DYNAMIC & Fa2/1 \\
1 & 0005.5ed3.1d1e (PC6) & DYNAMIC & Fa1/1 \\
1 & 0030.f223.90e0 (PC5) & DYNAMIC & Fa0/1 \\
\bottomrule
\end{tabular}
\caption{Q3 summary: ARP tables before/after (top) and the switch's final MAC forwarding
table (bottom). Before the pings the MAC table was empty.}
\end{table}

\paragraph{Observations.}
\begin{itemize}[leftmargin=1.6em]
  \item Both tables start empty and are filled \emph{only by traffic}. Neither is configured
  by hand; every entry is of type \texttt{dynamic}/\texttt{DYNAMIC}.
  \item The switch learns from the \emph{source} MAC of each frame arriving on a port. The
  first ARP request (a broadcast) teaches it the sender's port, and the unicast reply teaches
  it the target's port. Afterwards frames are forwarded only to the destination port instead
  of being flooded like on the hub.
  \item In the intermediate snapshot (after PC5's pings only) the entry for PC6 is missing,
  although PC6 had replied to PC5. It appeared in the next listing, after PC6 had itself
  sent pings. The final table has one entry per host, on the correct port.
  \item An ARP entry is learned in both directions by one exchange. The host that answers
  an ARP request also caches the requester's IP--MAC pair. That is why PC7's table already
  had two entries even though it had only pinged twice.
\end{itemize}

%--------------------------------------------------------------------------
\subsection{Q4: Joining LAN-A and LAN-B}

\subsubsection*{Topology and Addressing}
\shot{topo_q4.png}{Q4 topology: the hub of LAN-A and the switch of LAN-B joined by a
cross-over cable. Addresses are unchanged from Q2 and Q3.}

All six hosts keep their Q2/Q3 addresses (192.168.1.3--.8/24, MACs as in Tables~\ref{tab:lana} and~\ref{tab:lanb}).
The new link runs from a hub port to switch port \textbf{FastEthernet3/1}.

\subsubsection*{Procedure}
\begin{enumerate}[leftmargin=1.8em]
  \item Connected the hub (``LAN A'') to the switch (``LAN B'') with a \textbf{copper
  cross-over} cable. Hub-to-switch is a ``like device'' connection (both are
  interconnecting devices), so a cross-over cable is the correct choice. The switch end is
  Fa3/1.
  \item Pinged every LAN-A host from every LAN-B host and vice versa.
  \item Recorded \texttt{arp -a} on all six PCs and \texttt{show mac-address-table} on the
  switch.
\end{enumerate}

\subsubsection*{Results}
\shot{q4_arp_all.png}{\texttt{arp -a} on all six hosts after the cross-LAN pings. PC6 (top
left) and PC2--PC5, PC7 (right) each list the other \textbf{five} hosts.}
\shot{q4_mac_table.png}{Switch ``LAN B'' MAC table after the cross-LAN pings: six entries.
The three LAN-A hosts all appear on \textbf{Fa3/1}, the uplink to the hub.}

\begin{table}[H]
\centering
\begin{tabular}{lllll}
\toprule
\textbf{VLAN} & \textbf{MAC address} & \textbf{Host} & \textbf{Type} & \textbf{Port} \\
\midrule
1 & 0001.c70a.0697 & PC4 (LAN-A) & DYNAMIC & Fa3/1 \\
1 & 0002.16aa.b4dc & PC7 (LAN-B) & DYNAMIC & Fa2/1 \\
1 & 0005.5ed3.1d1e & PC6 (LAN-B) & DYNAMIC & Fa1/1 \\
1 & 0030.f223.90e0 & PC5 (LAN-B) & DYNAMIC & Fa0/1 \\
1 & 0090.2be8.b750 & PC3 (LAN-A) & DYNAMIC & Fa3/1 \\
1 & 00d0.ffa5.37ee & PC2 (LAN-A) & DYNAMIC & Fa3/1 \\
\bottomrule
\end{tabular}
\caption{MAC forwarding table of the LAN-B switch after Q4.}
\end{table}

\paragraph{Observations.}
\begin{itemize}[leftmargin=1.6em]
  \item Every cross-LAN ping succeeds without a router, because both LANs are in the same IP
  subnet and the hub--switch link merges them into one Layer-2 broadcast domain.
  \item Every ARP table now has five entries (all hosts except itself), versus two in Q3.
  \item Several MAC addresses mapping to \emph{one} switch port (Fa3/1) is the signature of
  a shared segment behind that port. The switch cannot tell that a hub is there. It only
  knows that those three MACs are reachable through Fa3/1. Traffic between two LAN-A hosts
  never needs the switch, and if it arrives on Fa3/1 for a MAC also on Fa3/1, the switch
  filters (drops) it rather than sending it back.
  \item The switch still separates collision domains: LAN-B hosts each have a private one,
  while the three LAN-A hosts and the switch's Fa3/1 share the hub's single collision domain.
\end{itemize}

%--------------------------------------------------------------------------
\subsection{Q5: Two Routed LANs (JU-Main and JU-SL) with Static Routes}

\subsubsection*{Topology and Addressing}
\shot{topo_q5.png}{Q5 topology: JU-Main and JU-SL, each a 2950 switch with three PCs behind
a 1841 router; the routers are joined by the 192.168.150.0/24 WAN link.}

\begin{table}[H]
\centering
\small
\resizebox{\textwidth}{!}{%
\begin{tabular}{lllll}
\toprule
\textbf{Device} & \textbf{Interface} & \textbf{IPv4 address / mask} & \textbf{MAC} & \textbf{Role} \\
\midrule
R1 (1841) & Fa0/0 & 192.168.148.1/24 & 0030.A3C5.8E01 & JU-Main gateway \\
R1 (1841) & Fa0/1 & 192.168.150.1/24 & 0030.A3C5.8E02 & WAN \\
R2 (1841) & Fa0/0 & 192.168.149.1/24 & 0010.11BE.0301 & JU-SL gateway \\
R2 (1841) & Fa0/1 & 192.168.150.2/24 & 0010.11BE.0302 & WAN \\
\midrule
PC8, PC9, PC10 & Fa0 & 192.168.148.2, .3, .4 /24 & -- & JU-Main hosts, gw 192.168.148.1 \\
PC12, PC13, PC14 & Fa0 & 192.168.149.2, .3, .4 /24 & -- & JU-SL hosts, gw 192.168.149.1 \\
\midrule
PC-LAB1-Switch & Fa0/1--0/4 & -- & base 0001.9770.0A01 & 2950-24, JU-Main \\
PC-LAB2-Switch & -- & -- & -- & 2950-24, JU-SL \\
\bottomrule
\end{tabular}}
\caption{Q5 addressing. Router addresses and MACs are from the device-info pop-ups
(Figures~\ref{fig:q5_r1_info.png} and~\ref{fig:q5_r2_info.png}).}
\label{tab:q5}
\end{table}

\begin{table}[H]
\centering
\begin{tabular}{llll}
\toprule
\textbf{Router} & \textbf{Destination network} & \textbf{Mask} & \textbf{Next hop} \\
\midrule
R1 & 192.168.149.0 & 255.255.255.0 & 192.168.150.2 (R2) \\
R2 & 192.168.148.0 & 255.255.255.0 & 192.168.150.1 (R1) \\
\bottomrule
\end{tabular}
\caption{Static routes. The 192.168.150.0/24 network and each router's own LAN are directly
connected and need no route.}
\label{tab:routes}
\end{table}

\subsubsection*{Procedure}
\begin{enumerate}[leftmargin=1.8em]
  \item \textbf{JU-Main:} placed a Cisco 2950-24 switch (hostname/label
  \texttt{PC-LAB1-Switch}), three PCs (PC8--PC10) and router R1 (1841). Cabled PCs and R1
  Fa0/0 to the switch with straight-through cables. Drew an ellipse labelled ``JU-Main''
  around it.
  \item \textbf{JU-SL:} the same with \texttt{PC-LAB2-Switch}, PC12--PC14 and R2.
  \item \textbf{WAN:} joined R1 Fa0/1 and R2 Fa0/1 with a cross-over cable (router to
  router is a like-device link). A FastEthernet point-to-point link was used instead of a
  serial link, because the 1841 has no serial interface without an add-on WIC module. The
  link plays the same WAN role, with its own subnet 192.168.150.0/24.
  \item \textbf{Router interfaces:} in each router's \emph{Config} tab, set the IP and mask
  of Fa0/0 and Fa0/1 and ticked \emph{Port Status: On} (Figures~\ref{fig:q5_r1_fa00.png},
  \ref{fig:q5_r1_fa01.png}). The equivalent IOS commands are:
\begin{lstlisting}
R1(config)# interface FastEthernet0/0
R1(config-if)# ip address 192.168.148.1 255.255.255.0
R1(config-if)# no shutdown
R1(config)# interface FastEthernet0/1
R1(config-if)# ip address 192.168.150.1 255.255.255.0
R1(config-if)# no shutdown
\end{lstlisting}
  \item \textbf{Hosts:} static IPs 192.168.148.2--.4 (JU-Main) and 192.168.149.2--.4
  (JU-SL), mask 255.255.255.0, and
  \emph{Default Gateway} = the router's Fa0/0 address (192.168.148.1 or 192.168.149.1).
  \item \textbf{Static routes:} \emph{Config $\to$ Routing $\to$ Static} on each router,
  with the entries of Table~\ref{tab:routes} (Figure~\ref{fig:q5_r1_static.png}). In IOS:
\begin{lstlisting}
R1(config)# ip route 192.168.149.0 255.255.255.0 192.168.150.2
R2(config)# ip route 192.168.148.0 255.255.255.0 192.168.150.1
\end{lstlisting}
  \item Tested from PC8: ping own gateway, the far end of the WAN, and the far LAN's gateway.
\end{enumerate}

\subsubsection*{Results: configuration}
\shot{q5_r1_fa00.png}{R1 \emph{Config $\to$ FastEthernet0/0}: port on, 192.168.148.1 /
255.255.255.0, MAC 0030.A3C5.8E01. The IOS pane below shows \texttt{show ip route} being run.}
\shot{q5_r1_fa01.png}{R1 \emph{Config $\to$ FastEthernet0/1}: 192.168.150.1 / 255.255.255.0.
The IOS pane shows ``192.168.150.0/24 is directly connected, FastEthernet0/1''.}
\shot{q5_r1_static.png}{R1 \emph{Config $\to$ Routing $\to$ Static}: the route
\texttt{192.168.149.0/24 via 192.168.150.2}.}
\shot{q5_r1_info.png}{Hover info for R1: Fa0/0 192.168.148.1/24 and Fa0/1 192.168.150.1/24,
both up.}
\shot{q5_r2_info.png}{Hover info for R2: Fa0/0 192.168.149.1/24 and Fa0/1 192.168.150.2/24,
both up.}
\shot{q5_switch_info.png}{Hover info for PC-LAB1-Switch (2950-24): Fa0/1--Fa0/4 up (three
PCs and R1), Fa0/5--Fa0/24 down, all in VLAN 1.}

\subsubsection*{Results: connectivity}
\shot{q5_ping_pc8_1.png}{PC8 pings its gateway 192.168.148.1 (TTL\,=\,255) and the far end
of the WAN link, 192.168.150.2 (TTL\,=\,254; first request timed out).}
\shot{q5_ping_pc8_2.png}{PC8 pings JU-SL's gateway 192.168.149.1: 4/4 replies, TTL\,=\,254.}

\paragraph{Observations.}
\begin{itemize}[leftmargin=1.6em]
  \item \textbf{Reachability proves both routes.} Echo \emph{requests} from PC8 to
  192.168.149.1 use R1's static route. The \emph{replies} back to 192.168.148.0/24 can only
  return through R2's static route to 192.168.148.0/24. Without that second route the
  requests would arrive and the replies would be dropped.
  \item \textbf{TTL counts the hops.} Router IOS sends replies with TTL 255. 192.168.148.1
  replies with 255, so zero routers lie in between. 192.168.150.2 and 192.168.149.1 reply
  with 254, because R1 forwarded the reply once.
  \item \textbf{The first lost packet} to 192.168.150.2 is again ARP. R1 had never resolved
  R2's Fa0/1 MAC and dropped the first packet while it did.
\end{itemize}

%--------------------------------------------------------------------------
\subsection{Q6: DHCP Servers in JU-Main and JU-SL}

\subsubsection*{Topology and Addressing}
\shot{topo_q6.png}{Q6 topology: Q5 plus one DHCP server in each LAN; hosts now obtain their
addresses dynamically.}

\begin{table}[H]
\centering
\small
\begin{tabular}{lll}
\toprule
\textbf{Setting} & \textbf{Server0 (JU-Main)} & \textbf{Server1 (JU-SL)} \\
\midrule
Server's own address (static) & 192.168.148.10/24 & 192.168.149.10/24 \\
Server's gateway / DNS & 192.168.148.1 / 8.8.8.8 & 192.168.149.1 / 8.8.8.8 \\
\midrule
DHCP pool name & \texttt{JU-Main} & \texttt{JU-SL} \\
Default gateway handed out & 192.168.148.1 & 192.168.149.1 \\
DNS server handed out & 8.8.8.8 & 8.8.8.8 \\
Start address & 192.168.148.100 & 192.168.149.100 \\
Subnet mask & 255.255.255.0 & 255.255.255.0 \\
Maximum users & 156 (.100--.255) & 156 (.100--.255) \\
\midrule
Leases observed & PC8 $\to$ 192.168.148.101 & PC13 $\to$ 192.168.149.101; \\
 & & 192.168.149.100 also leased \\
\bottomrule
\end{tabular}
\caption{DHCP configuration. Routers and WAN are unchanged from Table~\ref{tab:q5}.}
\end{table}

\subsubsection*{Procedure}
\begin{enumerate}[leftmargin=1.8em]
  \item Added a \texttt{Server-PT} to each LAN switch (Server0 in JU-Main, Server1 in JU-SL).
  \item Gave each server a \textbf{static} address (.10) outside the range it will hand out,
  with the router as its gateway (Figure~\ref{fig:q6_server_ip.png}).
  \item \emph{Services $\to$ DHCP}: turned the service \emph{On}, entered a new pool
  (\texttt{JU-Main} / \texttt{JU-SL}) with the values in the table and pressed \emph{Add}
  (Figure~\ref{fig:q6_dhcp_pools.png}).
  \item On every host: \emph{IP Configuration $\to$ DHCP}.
  \item Verified with \texttt{ipconfig} and a ping across both routers.
\end{enumerate}

\subsubsection*{Results}
\shot{q6_server_ip.png}{Static IP configuration of Server0 (192.168.148.10) and Server1
(192.168.149.10), each with its router as default gateway.}
\shot{q6_dhcp_pools.png}{DHCP service on both servers. Server0's pool \texttt{JU-Main}: gateway
192.168.148.1, start 192.168.148.100, 156 users. Server1's pool \texttt{JU-SL}: gateway
192.168.149.1, start 192.168.149.100.}
\shot{q6_pc13_dhcp.png}{PC13 (JU-SL) set to DHCP: ``DHCP request successful''. It received
192.168.149.101 / 255.255.255.0, gateway 192.168.149.1, DNS 8.8.8.8.}
\shot{q6_pc8_ipconfig.png}{\texttt{ipconfig} on PC8 (JU-Main): 192.168.148.101 /
255.255.255.0, default gateway 192.168.148.1, all from Server0.}
\shot{q6_pc8_ping.png}{PC8 (DHCP address) pings a DHCP-configured host in JU-SL,
192.168.149.100: 3 of 4 replies, TTL\,=\,126.}

\paragraph{Observations.}
\begin{itemize}[leftmargin=1.6em]
  \item The address, mask \emph{and} default gateway all came from DHCP. The gateway is
  what lets the DHCP hosts keep using Q5's static routes; the cross-LAN ping succeeds.
  \item TTL\,=\,126 on the reply is the PC default 128 minus two router hops (R2 and R1),
  confirming the packet crossed both routers. The first request was lost to ARP resolution
  at R2 for the new host's address.
  \item A DHCP DISCOVER is a broadcast and routers do not forward broadcasts. That is why
  each LAN needs its own DHCP server, as the assignment specifies. The alternative is a
  single server plus \texttt{ip helper-address} on each router's LAN interface.
  \item The pool starts at .100, so it never collides with the statically addressed router
  (.1) and server (.10). The hosts move from their Q5 static addresses (.2--.4) to leases
  from .100 onwards.
  \item DNS 8.8.8.8 is a placeholder required by the form; no DNS is used in this question.
\end{itemize}

%--------------------------------------------------------------------------
\subsection{Q7: The CSE LAN with DHCP, DNS, Web and FTP Servers}

\subsubsection*{Topology and Addressing}
\shot[0.92\textwidth]{topo_q7.png}{Q7 topology: the CSE LAN. Three DHCP clients and four
statically addressed servers on one 2950 switch.}

\begin{table}[H]
\centering
\small
\resizebox{\textwidth}{!}{%
\begin{tabular}{lllll}
\toprule
\textbf{Device} & \textbf{IPv4 address} & \textbf{How assigned} & \textbf{DNS server} & \textbf{DNS name} \\
\midrule
DHCP Server & 192.168.151.10/24 & static & 192.168.151.13 & -- \\
WebServer & 192.168.151.11/24 & static & 192.168.151.13 & \texttt{www.cse.myuniv.edu} \\
FTP Server & 192.168.151.12/24 & static & 192.168.151.13 & \texttt{ftp.cse.myuniv.edu} \\
DNS Server & 192.168.151.13/24 & static & 192.168.151.13 & -- \\
\midrule
CSE-PC1 & 192.168.151.100/24 & DHCP & 192.168.151.13 (via DHCP) & -- \\
CSE-PC2, CSE-PC3 & from pool (.100 onwards) & DHCP & 192.168.151.13 (via DHCP) & -- \\
\midrule
CSE-Switch & -- & -- & -- & Cisco 2950-24 \\
\bottomrule
\end{tabular}}
\caption{Q7 addressing. The LAN has no router, so the default gateway is 0.0.0.0 everywhere.}
\end{table}

\subsubsection*{Procedure}
\begin{enumerate}[leftmargin=1.8em]
  \item Placed a 2950-24 switch \texttt{CSE-Switch}, three PCs (CSE-PC1--3) and four
  \texttt{Server-PT}s (DHCP Server, WebServer, FTP Server, DNS Server), all on
  straight-through cables.
  \item \textbf{Static server addresses} 192.168.151.10--.13/24, each with DNS server
  192.168.151.13 (Figure~\ref{fig:q7_server_ips.png}).
  \item \textbf{DHCP Server:} \emph{Services $\to$ DHCP} On; pool \texttt{serverPool}
  with start 192.168.151.100, mask 255.255.255.0, \textbf{DNS server 192.168.151.13}, gateway
  0.0.0.0, maximum users 50 (Figure~\ref{fig:q7_dhcp_pool.png}). Handing out the DNS address
  through DHCP is what ``configures the individual hosts to use the local DNS server''.
  \item \textbf{DNS Server:} \emph{Services $\to$ DNS} On; added A records
  (Figures~\ref{fig:q7_dns_record1.png}, \ref{fig:q7_dns_records.png}):
  \begin{center}
  \begin{tabular}{lll}
  \toprule
  \textbf{Name} & \textbf{Type} & \textbf{Address} \\
  \midrule
  \texttt{www.cse.myuniv.edu} & A & 192.168.151.11 \\
  \texttt{ftp.cse.myuniv.edu} & A & 192.168.151.12 \\
  \bottomrule
  \end{tabular}
  \end{center}
  A first draft had a single record for the bare domain \texttt{cse.myuniv.edu}. It was
  replaced by one host name per server under the domain.
  \item \textbf{WebServer:} \emph{Services $\to$ HTTP}; edited \texttt{index.html} to a CSE
  welcome page (Figure~\ref{fig:q7_web_index.png}).
  \item \textbf{FTP Server:} \emph{Services $\to$ FTP} On; added user \texttt{sajjad}
  (permissions Read, Write, List) and \texttt{admin} (all permissions) alongside the default
  \texttt{cisco} (Figure~\ref{fig:q7_ftp_users.png}).
  \item \textbf{Hosts:} \emph{IP Configuration $\to$ DHCP}.
  \item Tested name resolution with \texttt{ping} by name, \texttt{http://www.cse.myuniv.edu}
  in the PC's Web Browser, and \texttt{ftp ftp.cse.myuniv.edu} from the command prompt.
\end{enumerate}

\subsubsection*{Results: configuration}
\shot{q7_server_ips.png}{Static IP configuration of the four servers: DHCP .10, WebServer .11,
FTP .12, DNS .13. All point to 192.168.151.13 for DNS.}
\shot{q7_dhcp_pool.png}{DHCP Server pool: DNS 192.168.151.13, start 192.168.151.100 /
255.255.255.0, maximum users being set to 50.}
\shot{q7_dns_record1.png}{DNS Server, first draft: one A record \texttt{cse.myuniv.edu}
$\to$ 192.168.151.11.}
\shot{q7_dns_records.png}{DNS Server, final records: \texttt{www.cse.myuniv.edu} $\to$
192.168.151.11 and \texttt{ftp.cse.myuniv.edu} $\to$ 192.168.151.12.}
\shot{q7_web_index.png}{WebServer HTTP service: the edited \texttt{index.html}.}
\shot{q7_ftp_users.png}{FTP Server users: \texttt{cisco} (RWDNL), \texttt{sajjad} (RWL),
\texttt{admin} (RWDNL). R/W/D/N/L = Read, Write, Delete, reName, List.}
\shot{q7_pc1_dhcp.png}{CSE-PC1 on DHCP: 192.168.151.100 / 255.255.255.0, gateway 0.0.0.0,
DNS server 192.168.151.13, all received from the DHCP Server.}

\subsubsection*{Results: services}
\shot{q7_ping_names.png}{CSE-PC1 pings by name: \texttt{www.cse.myuniv.edu} resolves to
192.168.151.11 and \texttt{ftp.cse.myuniv.edu} to 192.168.151.12. 0\% loss on both.}
\shot{q7_browser.png}{CSE-PC1's Web Browser at \texttt{http://www.cse.myuniv.edu}: the page
served by WebServer.}
\shot{q7_ftp_session.png}{FTP session from CSE-PC1: \texttt{ftp ftp.cse.myuniv.edu},
\texttt{220 Welcome}, user \texttt{sajjad}, \texttt{230 Logged in}, then \texttt{dir} lists
the server's \texttt{/ftp} directory. (Packet Tracer's client does not support \texttt{ls}.)}

\paragraph{Observations.}
\begin{itemize}[leftmargin=1.6em]
  \item One chain links all four services. DHCP gives CSE-PC1 an address \emph{and} the DNS
  server's address. DNS turns \texttt{www}/\texttt{ftp.cse.myuniv.edu} into .11/.12. HTTP
  and FTP then connect to those addresses. The \texttt{Pinging 192.168.151.11} line shows the
  resolved address that \texttt{ping} actually used.
  \item Servers are static, so their DNS records never go stale. Hosts are dynamic, and
  nothing needs to resolve \emph{them} by name.
  \item The FTP exchange shows the protocol's numeric replies: 220 (service ready), 331
  (user OK, need password), 230 (logged in).
  \item No router or gateway is needed, because every device is on 192.168.151.0/24.
\end{itemize}

%==========================================================================
\section{Analysis}

\subsection{Hub vs.\ switch, from the tables}
Q2--Q4 show the difference between Layer 1 and Layer 2 directly. The hub has no table at
all; it repeats every bit. The switch starts with an empty MAC table and builds it purely
from source addresses (Q3). In Q4 three MACs share port Fa3/1, which is how a switch sees
a hub (or another switch) behind one of its ports. ARP is the bridge between Layers 2 and 3.
Its tables fill only on demand, and a single request/reply exchange teaches \emph{both}
hosts.

\subsection{What a router changes}
Q4 joins two LANs \emph{without} a router, and it works only because they share one subnet
and become one broadcast domain. Q5 joins two LANs in \emph{different} subnets, which needs
(a) a router interface in each subnet, (b) a default gateway on each host, and (c) a route
on each router for every non-attached network, in \emph{both} directions. The TTL values make
the router hops visible: 128 between PCs on one LAN, 255/254 from the routers, 126 across
two routers.

\subsection{From manual to automatic configuration}
Q5 configured every host by hand. Q6 replaced that with DHCP: one pool per LAN delivers the
address, mask and gateway. Q7 adds DNS: hosts learn the DNS server from DHCP, and users name
services instead of addresses. Each layer of automation depends on the one below: DHCP
needs Layer-2 broadcast reachability, DNS needs a correct DHCP option, and HTTP/FTP need
DNS.

\subsection{The ``first ping lost'' pattern}
Q2, Q5 and Q6 each lost exactly one request the first time a new next-hop was used. Every
case was an ARP miss. Packet Tracer (like many real IP stacks) drops the packet that
triggers ARP resolution instead of queuing it. It is not a configuration fault, and a
repeated ping always showed 0\% loss.

%==========================================================================
\section{Comments}

\paragraph{Evaluation of the Lab Assignment.} The assignment builds up one layer at a
time, and each question reuses the previous one. LAN-A and LAN-B from Q2/Q3 are joined in
Q4, and the routed network of Q5 gains DHCP in Q6. Recording the ARP and MAC tables before
and after traffic made address learning concrete in a way the earlier socket-based
assignments could not.

\paragraph{Difficulty.} The difficulty was mostly in Packet Tracer details. The router
model named in the problem (``1818'') does not exist, so the 1841 was used. The 1841 has no
serial port without a WIC module, so the WAN link was made with FastEthernet. The IOS
syntax is \texttt{show mac address-table} (a mistyped \texttt{show mac-address table} is
rejected, Figure~\ref{fig:q3_mac_before.png}). The PC's FTP client accepts \texttt{dir} but
not \texttt{ls}. The one-lost-ping-on-ARP behaviour also looks like a fault until its cause
is understood.

\paragraph{Suggested Improvements.} Correct the router model (1841 or 1941), and say
whether the WAN link should be serial (which would teach clock rate and DCE/DTE) or
Ethernet. In Q7, ask for a router as well, so that the CSE LAN can be connected to the
JU-Main/JU-SL network and DNS names resolved across subnets.

\end{document}
